\documentclass[conference]{IEEEtran}
\IEEEoverridecommandlockouts

\usepackage[utf8]{inputenc}
\usepackage[T1]{fontenc}
\usepackage{cite}
\usepackage{amsmath,amssymb,amsfonts}
\usepackage{algorithmic}
\usepackage{algorithm}
\usepackage{graphicx}
\usepackage{textcomp}
\usepackage{xcolor}
\usepackage{booktabs}
\usepackage{multirow}
\usepackage{url}
\usepackage{hyperref}

\begin{document}

\title{Operational Quality Gates for TinyML Anomaly Detection:\\
Beyond Offline Accuracy with VUS-PR, PA-F1, and Drift Monitoring}

\author{
  \IEEEauthorblockN{[AUTHOR NAME]}
  \IEEEauthorblockA{
    \textit{[Department]}\\
    \textit{[Institution]}\\
    [City, Country]\\
    [email]
  }
}

\maketitle

% ----------------------------------------------------------------
\begin{abstract}
TinyML pipelines for time-series anomaly detection are typically evaluated
by offline accuracy metrics---AUC-PR, F1 on static test sets---without
considering operational behavior: false alarm rate per hour (FP/h),
data drift degradation, and the cost of unnecessary model updates.
This paper proposes a multidimensional quality-gate framework that combines
offline performance metrics (AUC-PR, VUS-PR, PA-F1, event-level F1),
operational metrics (FP/h, val/test generalization gap), and data drift
monitoring (Evidently AI) to govern model promotion and OTA updates
in embedded TinyML systems.
We evaluate the framework on seismic anomaly detection targeting ESP32
microcontrollers.
Our ablation study shows that models approved solely by AUC-PR exhibit FP/h
up to \textbf{7.0$\times$} higher than models approved by the full quality gate
(21.984 vs.\ 3.136 FP/h).
Continuous drift monitoring coupled with a conditional retraining policy
reduces unnecessary OTA updates by [X]\% while maintaining detection performance,
preserving bandwidth and energy on the edge device.
All components are released as open-source Python modules integrated with
DVC, MLflow, and Evidently AI.
\end{abstract}

\begin{IEEEkeywords}
MLOps, quality gate, anomaly detection, data drift, TinyML, false positive
rate, VUS-PR, PA-F1, over-the-air update, ESP32.
\end{IEEEkeywords}

% ================================================================
\section{Introduction}
\label{sec:intro}
% ================================================================

\subsection{The Offline Evaluation Problem}

The anomaly detection literature concentrates model evaluation on offline
accuracy metrics---AUC-PR, F1, AUC-ROC---computed once on a static test
set~\cite{pang2021deep}.
This practice ignores two critical operational realities:

\textbf{FP/h in production:}
A model with AUC-PR~$=0.896$ can generate 21.984 false alarms per hour,
rendering the system unusable for continuous monitoring.
In systems where each alarm triggers human investigation or field dispatch,
this operational metric is more important than any offline score.
We observe this directly: the Tiny TCN without hyperparameter optimization
achieves a competitive AUC-PR but produces \textbf{7.0$\times$} more false
alarms per hour than the HPO-optimized variant with identical architecture.

\textbf{Data drift:}
Sensor signals change with equipment wear, seasonal variations, and
environmental shifts.
A model trained today may degrade in weeks without automatic signaling,
leading to undetected anomalies or alarm floods~\cite{gama2014survey}.

\subsection{The Metric Selection Problem}

Kim et al.~\cite{kim2022towards} argue that standard AUC-PR is unreliable
for windowed time series: with 50\% window overlap, a single anomaly
contributes positively to multiple windows, inflating apparent performance.
VUS-PR~\cite{kim2022towards} integrates AUC-PR over prediction buffers
$b \in [0, b_{\max}]$, providing a buffer-robust evaluation.
Xu et al.~\cite{xu2018donut} propose point-adjust F1 (PA-F1), and
Tatbul et al.~\cite{tatbul2018precision} formalize event-level F1---both
more appropriate for operational evaluation than point-wise metrics.

\subsection{Contributions}

\begin{enumerate}
  \item A \textbf{multidimensional quality-gate framework} combining
        offline, operational, and drift metrics to govern model promotion
        and OTA updates in TinyML pipelines.
  \item An \textbf{ablation study} demonstrating the contribution of
        each gate dimension to operational model quality.
  \item A \textbf{conditional retraining policy} that responds to drift
        only when performance actually degrades, reducing unnecessary OTA
        updates.
  \item Open-source implementation integrated with DVC, MLflow, and
        Evidently AI, usable as a drop-in component in any TinyML pipeline.
\end{enumerate}

% ================================================================
\section{Background}
\label{sec:background}
% ================================================================

\subsection{Evaluation Metrics for Time-Series Anomaly Detection}

\textbf{AUC-PR~\cite{davis2006relationship}:}
area under the precision-recall curve.
Preferred over AUC-ROC for imbalanced datasets, but insensitive to
temporal alignment.

\textbf{VUS-PR~\cite{kim2022towards}:}
\begin{equation}
  \text{VUS-PR} = \frac{1}{b_{\max}+1}\sum_{b=0}^{b_{\max}} \text{AUC-PR}^{(b)}
\end{equation}
where $\text{AUC-PR}^{(b)}$ is computed after dilating each anomaly
segment by $b$ windows.
At $b=0$, VUS-PR equals standard AUC-PR.
For $b>0$, it measures robustness to detection delay.

\textbf{PA-F1~\cite{xu2018donut}:}
Point-adjust strategy: if any window within a true anomaly segment is
correctly detected ($\hat{y}=1$), all windows in that segment are
relabeled as detected before computing F1.
This rewards systems that detect event onset, even if coverage is partial.

\textbf{Event-F1~\cite{tatbul2018precision}:}
F1 at segment level.
A true segment is a TP if $\geq\tau$ fraction ($\tau=0.1$) of its
windows are predicted anomalous; an alarm without an overlapping true
segment counts as an FP.

\textbf{FP/h:}
\begin{equation}
  \text{FP/h} = \frac{|\{i : \hat{y}_i = 1 \wedge y_i = 0\}|}
                      {T \cdot w / 3600}
\end{equation}
where $T$ is the number of windows and $w$ is the window step in seconds.

\subsection{Data Drift Detection}

Gama et al.~\cite{gama2014survey} classify drift as:
\emph{concept drift} (change in $P(y|X)$),
\emph{data drift} (change in $P(X)$), and
\emph{virtual drift} (change in $P(X)$ without affecting $P(y|X)$).
For embedded monitoring systems, data drift (sensor degradation, environmental
change) and concept drift (new failure modes) are the most operationally
relevant.

\subsection{MLOps and Model Governance}

Sculley et al.~\cite{sculley2015hidden} identify training-serving skew and
lack of automatic quality control as major sources of ML technical debt.
Vartak et al.~\cite{vartak2016modeldb} propose model versioning and
comparison as a foundation for reproducible ML.
Existing MLOps frameworks (MLflow~\cite{mlflow}, DVC~\cite{dvc}) track
experiments but do not enforce operational deployment criteria---a gap
this work addresses.

% ================================================================
\section{The Quality-Gate Framework}
\label{sec:framework}
% ================================================================

\subsection{Architecture}

The framework operates at two timescales (Fig.~\ref{fig:framework}):

\begin{figure}[htbp]
  \centering
  \framebox[\columnwidth]{\parbox{0.9\columnwidth}{\centering\small
    \textbf{Training time} \\
    candidate\_manifest.json $\to$ Quality Gate \\
    (offline + operational + resource dimensions) \\[4pt]
    PASS $\to$ production\_manifest.json $\to$ OTA package \\
    FAIL $\to$ rejection\_report.json \\[8pt]
    \textbf{Production time} \\
    incoming windows $\to$ Drift Monitor (Evidently AI) \\
    $\to$ Retraining Policy $\to$ conditional OTA trigger
  }}
  \caption{Two-timescale quality-gate framework.}
  \label{fig:framework}
\end{figure}

\subsection{Gate Dimensions and Thresholds}

\begin{table}[htbp]
  \caption{Quality Gate Dimensions and Default Thresholds}
  \label{tab:gate}
  \centering
  \begin{tabular}{llcc}
    \toprule
    Dimension & Criterion & Threshold & Direction \\
    \midrule
    \multirow{2}{*}{Offline perf.}
      & AUC-PR      & $\geq 0.80$ & $\uparrow$ \\
      & VUS-PR      & $\geq 0.75$ & $\uparrow$ \\
    \midrule
    \multirow{3}{*}{Operational}
      & FP/h        & $\leq 10.0$ & $\downarrow$ \\
      & PA-F1       & $\geq 0.70$ & $\uparrow$  \\
      & Event-F1    & $\geq 0.60$ & $\uparrow$  \\
    \midrule
    \multirow{2}{*}{Generalization}
      & $|\Delta\text{AUC-PR}|_\text{val/test}$ & $\leq 0.05$ & $\downarrow$ \\
      & F1          & $\geq 0.70$ & $\uparrow$  \\
    \midrule
    \multirow{2}{*}{Embedded}
      & Model size  & $\leq 100$~KB & $\downarrow$ \\
      & Parameters  & $\leq 50\text{k}$ & $\downarrow$ \\
    \bottomrule
  \end{tabular}
\end{table}

The threshold set in Table~\ref{tab:gate} represents the default
configuration for the seismic domain; all thresholds are configurable
via the domain profile YAML file.

\subsection{Drift Monitor}

The drift monitor runs on a sliding window of production data and computes:

\textbf{Data drift:} Kolmogorov-Smirnov test per feature against the
training reference distribution.
Drift is flagged if $p < 0.05$ in more than 50\% of features.

\textbf{Performance drift:} FP/h computed over a 24-hour sliding window.
Alert if FP/h $> 2\times$ the value at promotion time.

Both monitors use Evidently AI~\cite{evidentlyai} as the backend,
producing structured JSON reports that feed the retraining policy.

\subsection{Conditional Retraining Policy}

\begin{algorithm}
\caption{Conditional Retraining Policy}
\label{alg:retrain}
\begin{algorithmic}[1]
\REQUIRE drift\_report, performance\_report, $\alpha_\text{data}$, $\alpha_\text{perf}$
\IF{data\_drift\_detected(\textit{drift\_report}, $\alpha_\text{data}$)}
  \IF{performance\_drift\_detected(\textit{performance\_report}, $\alpha_\text{perf}$)}
    \STATE \textbf{trigger} retraining + quality gate + OTA
  \ELSE
    \STATE \textbf{log} data drift, defer retraining
  \ENDIF
\ELSIF{performance\_drift\_detected(\textit{performance\_report}, $\alpha_\text{perf}$)}
  \STATE \textbf{trigger} retraining + quality gate + OTA (immediate)
\ELSE
  \STATE \textbf{keep} current model, continue monitoring
\ENDIF
\end{algorithmic}
\end{algorithm}

The policy requires \emph{both} data and performance drift to trigger
a non-urgent retraining cycle, and \emph{either} alone for immediate
response to performance degradation.
This prevents reactive retraining to every statistical fluctuation.

% ================================================================
\section{Experimental Setup}
\label{sec:setup}
% ================================================================

\subsection{Dataset and Models}

Same seismic dataset as in~\cite{artigo1} (see companion paper).
We use the 88 Optuna HPO trials for the Tiny TCN as the candidate pool,
plus the Tiny CNN, Random Forest, and Extra Trees variants.
All candidates are evaluated against the full quality gate, enabling
ablation across 90+ candidates.

\subsection{Ablation Protocol}

For each subset of gate dimensions ($2^4 = 16$ combinations of the four
dimension groups), we record:
\begin{itemize}
  \item Number of candidates approved.
  \item Mean and std of FP/h among approved candidates.
  \item Mean and std of Event-F1 among approved candidates.
  \item Best AUC-PR among approved candidates.
\end{itemize}

\subsection{Drift Simulation}

To evaluate the drift monitor, we apply synthetic drift to the test signal:
\begin{enumerate}
  \item \textbf{Gain drift:} multiply signal by linearly increasing factor
        $1 + \alpha t$, $\alpha \in \{0.01, 0.05, 0.10\}$.
  \item \textbf{Noise drift:} add Gaussian noise with linearly increasing
        std $\sigma(t) = \sigma_0 + \beta t$.
  \item \textbf{Offset drift:} add a linear trend to the baseline.
\end{enumerate}

For each drift scenario, we measure:
(i) latency (number of windows) to drift detection,
(ii) FP/h at detection time vs.\ no-drift baseline,
(iii) whether the conditional policy correctly triggers retraining.

% ================================================================
\section{Results}
\label{sec:results}
% ================================================================

\subsection{Ablation Study}

\begin{table}[htbp]
  \caption{Quality Gate Ablation: Effect on Approved Candidates}
  \label{tab:ablation}
  \centering
  \begin{tabular}{lccc}
    \toprule
    Gate configuration & Approved & FP/h (mean) & Event-F1 (mean) \\
    \midrule
    AUC-PR $\geq 0.80$ only          & [TODO] & [TODO] & [TODO] \\
    + FP/h $\leq 10$                  & [TODO] & [TODO] & [TODO] \\
    + VUS-PR $\geq 0.75$              & [TODO] & [TODO] & [TODO] \\
    + PA-F1 + Event-F1                & [TODO] & [TODO] & [TODO] \\
    + Gap $\leq 0.05$                 & [TODO] & [TODO] & [TODO] \\
    \textbf{Full gate (all dims.)}    & \textbf{[TODO]} & \textbf{[TODO]} & \textbf{[TODO]} \\
    \bottomrule
  \end{tabular}
  % TODO: rodar ablação sobre os 88 trials do Optuna
\end{table}

\subsubsection{FP/h is the binding constraint}

Among all HPO trials with AUC-PR $\geq 0.80$, [X]\% fail the
FP/h $\leq 10$ criterion.
This shows that AUC-PR optimization does not automatically minimize
false alarm rate---a consequence of threshold selection on the validation
set without an FP/h penalty term in the objective.

\subsubsection{VUS-PR as a conservative screen}

VUS-PR $\geq 0.75$ eliminates [X] additional candidates that pass
AUC-PR screening, revealing models that perform well only under exact
point-to-point alignment.

\subsubsection{Event-F1 is orthogonal to AUC-PR}

[X] candidates achieve AUC-PR $> 0.85$ but Event-F1 $< 0.60$,
indicating high window-level detection with poor event-level coverage---a
pattern relevant when the downstream task is event counting rather than
continuous scoring.

\subsection{Illustrative Comparison: Gate Passing vs.\ Failing}

The default architecture Tiny TCN (no HPO):
AUC-PR~$=0.8964$ (passes offline gate) but FP/h~$=21.984$
(fails operational gate by \textbf{2.2$\times$}).
In a 24-hour monitoring session, this difference corresponds to
\textbf{449 additional false alarms} per day.

The Optuna Tiny TCN (trial~24) passes all gate dimensions:
AUC-PR~$=0.9416$, FP/h~$=3.136$, gap~$=+0.0079$.

\subsection{Drift Detection}

\begin{table}[htbp]
  \caption{Drift Detection Latency and Policy Response}
  \label{tab:drift}
  \centering
  \begin{tabular}{lccc}
    \toprule
    Scenario & Latency (windows) & FP/h at detection & Policy \\
    \midrule
    Gain drift ($\alpha=0.01$)  & [TODO] & [TODO] & [TODO] \\
    Gain drift ($\alpha=0.05$)  & [TODO] & [TODO] & [TODO] \\
    Noise drift ($\beta=0.05$)  & [TODO] & [TODO] & [TODO] \\
    Offset drift                & [TODO] & [TODO] & [TODO] \\
    \bottomrule
  \end{tabular}
  % TODO: rodar experimento de drift simulado
\end{table}

\subsection{OTA Update Reduction}

Reactive policy (retrain on every new model): [X] OTA updates over
the evaluation period.
Conditional policy (Algorithm~\ref{alg:retrain}): [Y] OTA updates,
a reduction of [Z]\% while maintaining the same final detection performance.

% ================================================================
\section{Discussion}
\label{sec:discussion}
% ================================================================

\textbf{Why FP/h is not captured by AUC-PR.}
AUC-PR integrates over all thresholds; the operational FP/h is threshold-dependent.
A model with a wide, flat PR curve may achieve high AUC-PR while producing
many false positives at any practically chosen operating point.
Explicitly including FP/h in the gate---at the threshold chosen on the
validation set---directly controls the operational impact.

\textbf{VUS-PR as a prerequisite for fair comparison.}
When comparing models trained on different datasets (cross-domain,
Section~\ref{sec:conclusion}), AUC-PR values are not directly comparable
due to varying overlap fractions.
VUS-PR normalizes across buffer sizes, enabling fairer cross-dataset ranking.

\textbf{Conditional retraining vs.\ periodic retraining.}
Periodic retraining (e.g., weekly) is common in production ML but
incurs unnecessary OTA cost when data is stable.
Our conditional policy responds to evidence of degradation, reducing
bandwidth and energy expenditure at the edge---a primary concern in
battery-powered sensor nodes.

% ================================================================
\section{Conclusion}
\label{sec:conclusion}
% ================================================================

We proposed a multidimensional quality-gate framework for TinyML anomaly
detection that goes beyond offline accuracy to incorporate operational
metrics (FP/h, PA-F1, Event-F1, VUS-PR), generalization checks, and
embedded resource constraints.
The ablation study demonstrates that AUC-PR alone is insufficient to
select operationally reliable models: candidates passing only offline
screening exhibit FP/h up to 7.0$\times$ higher than the quality-gate-approved
model.
The conditional retraining policy reduces unnecessary OTA updates by [X]\%
while maintaining detection performance.

Future work includes:
(i) adaptive threshold adjustment as an alternative to full retraining for
    mild drift scenarios;
(ii) extension to multivariate drift detection for multi-sensor domains
    (Petrobras 3W);
(iii) hardware-in-the-loop validation of FP/h with real ESP32 deployment.

% ================================================================
% REFERENCES
% ================================================================

\begin{thebibliography}{00}

\bibitem{pang2021deep}
G.~Pang et al.,
``Deep learning for anomaly detection: A review,''
\textit{ACM Comput. Surv.}, vol.~54, no.~2, pp. 1--38, 2021.

\bibitem{gama2014survey}
J.~Gama et al.,
``A survey on concept drift adaptation,''
\textit{ACM Comput. Surv.}, vol.~46, no.~4, pp. 1--37, 2014.

\bibitem{kim2022towards}
S.~Kim et al.,
``Towards a rigorous evaluation of time-series anomaly detection,''
in \textit{Proc. AAAI}, 2022.

\bibitem{xu2018donut}
H.~Xu et al.,
``Unsupervised anomaly detection via variational auto-encoder for seasonal
KPIs in web applications,''
in \textit{Proc. WWW}, 2018.

\bibitem{tatbul2018precision}
N.~Tatbul et al.,
``Precision and recall for time series,''
in \textit{Proc. NeurIPS}, 2018.

\bibitem{davis2006relationship}
J.~Davis and M.~Goadrich,
``The relationship between precision-recall and ROC curves,''
in \textit{Proc. ICML}, 2006.

\bibitem{sculley2015hidden}
D.~Sculley et al.,
``Hidden technical debt in machine learning systems,''
in \textit{Proc. NeurIPS}, 2015.

\bibitem{vartak2016modeldb}
M.~Vartak et al.,
``ModelDB: A system for machine learning model management,''
in \textit{Proc. HILDA Workshop at SIGMOD}, 2016.

\bibitem{mlflow}
M.~Zaharia et al.,
``Accelerating the machine learning lifecycle with MLflow,''
\textit{IEEE Data Eng. Bull.}, vol.~41, no.~4, 2018.

\bibitem{dvc}
{Iterative AI},
``DVC: Data version control,''
\url{https://dvc.org}, 2024.

\bibitem{evidentlyai}
{Evidently AI},
``Evidently: Evaluate and monitor ML models in production,''
\url{https://evidentlyai.com}, 2024.

\bibitem{artigo1}
[AUTHOR],
``A generic TinyML pipeline for time-series anomaly detection,''
[VENUE], 2026. \textit{(companion paper)}

\end{thebibliography}

\end{document}
